
\section{Related Work}

\subsection{Static Sparse Attention}

\textbf{Fixed-pattern sparsity} restricts attention to a predetermined subset of 
key positions:

\begin{itemize}
\item \textbf{Longformer} (Beltagy et al., 2020): Sliding window + global tokens.
\item \textbf{BigBird} (Zaheer et al., 2020): Sliding window + random + global.
\item \textbf{Sparse Transformer} (Child et al., 2019): Strided + local patterns.
\item \textbf{BlockBERT} (Qiu et al., 2020): Block-sparse pattern.

\end{itemize}
\textbf{Limitation.} Patterns are fixed before training and cannot adapt to 
content or head heterogeneity.

\subsection{Dynamic Sparse Attention}

\textbf{Content-dependent sparsity} selects keys based on query-key similarity:

\begin{itemize}
\item \textbf{Routing Transformer} (Roy et al., 2021): K-means clustering of queries
  and keys.
\item \textbf{Reformer} (Kitaev et al., 2020): Locality-sensitive hashing.
\item \textbf{Sinkhorn Attention} (Tay et al., 2020): Sorted bucket attention.
\item \textbf{Top-K Attention} (Gupta et al., 2021): Hard top-$k$ over scores.

\end{itemize}
\textbf{Limitation.} Most dynamic methods use \textbf{uniform} sparsity across heads, 
or use hard top-$k$ with zero gradient.

\subsection{Learned Sparsity}

\textbf{Learnable patterns} adapt the sparsity structure during training:

\begin{itemize}
\item \textbf{Learned Sparsity Patterns} (Tay et al., 2021): Differentiable stride and
  dilation.
\item \textbf{SangerNet} (Yuan et al., 2021): Learned sparse masks.
\item \textbf{SparseBERT} (Xia et al., 2022): Distilled sparse patterns.

\end{itemize}
\textbf{ASA differs} by learning \textbf{per-head} budgets via a \textbf{scalar} parameter 
per head, avoiding the complexity of learnable masks.

\subsection{Differentiable Top-$k$}

\textbf{Hard top-$k$} has zero gradient. Several relaxations exist:

\begin{itemize}
\item \textbf{Soft top-$k$} (Augustin et al., 2022): Perturbation-based.
\item \textbf{Gumbel top-$k$} (Kool et al., 2019): Stochastic.
\item \textbf{Sigmoid soft mask} (our approach): Deterministic, based on threshold.

\end{itemize}
\textbf{ASA uses the sigmoid soft mask}, which is simpler than existing relaxations 
and works well empirically.

\subsection{Head Importance and Sparsity}

\textbf{Head pruning} studies which heads matter:

\begin{itemize}
\item Voita et al. (2019): Most heads can be pruned with minimal loss.
\item Michel et al. (2019): A single head suffices for many tasks.
\item Clark et al. (2019): Attention patterns vary significantly across heads.

\end{itemize}
\textbf{ASA exploits this} by allowing each head to learn its own budget, 
effectively performing soft head pruning.

\subsection{Position of Our Work}

\begin{table}[H]
\centering
\small
\begin{tabular}{lll}
\toprule
\textbf{Axis} & \textbf{Prior work} & \textbf{ASA} \\
\midrule
Sparsity pattern \& Fixed or uniform & \textbf{Per-head learned} \\
Top-$k$ selection \& Hard (non-diff) & \textbf{Soft (differentiable)} \\
Head budget \& Uniform & \textbf{Heterogeneous} \\
Training \& From scratch & \textbf{Pretrained + LoRA} \\
\bottomrule
\end{tabular}
\end{table}

\textbf{Key distinction.} ASA is the first method to learn **per-head sparsity 
budgets\textbf{ via a }single scalar per head**, without auxiliary losses or 
special training procedures.

\subsection{Comparison Table}

\begin{lstlisting}

---
\end{lstlisting}